\section{The slice estimate and the supremum bootstrap}
\label{boot:section}

Throughout this section, $n=N+1\geq 5$, and the local theorem is
assumed in dimensions smaller than $n$.  Recall that
$p_n=2(N/n)^n$ and that $M_n$ is the supremum of $d(x,X)$ over
all singular boundary-zero klt germs of dimension $n$.
The smooth-value upper bound gives $M_n\leq 1$; no bound by $p_n$
is assumed.  We first work with a fixed singular cone $o\in C=\Spec S$
obtained from Proposition~\ref{cone:reduction}, such that
$d(o,C)\geq p_n$.  Its minimizing Reeb vector $\xi$ is normalized by
$A_C(\wt_\xi)=n$.

Choose a sufficiently fast grading from
Lemma~\ref{cone:approximation} in the large-stabilizer range,
and take $m=m_{\max}\geq r/N$.  All constructions in
Proposition~\ref{filt:structure} then apply.  We shall use, in particular,
\begin{equation}
 \label{boot:algebra-data}
 \bar S=\gr_T S,
 \qquad D=\bar S/(\bar v),
 \qquad G=\gr_W S=D[\mathtt v],
 \qquad \Hilb_G(z)=(1-z)^{-n}.
\end{equation}
Here $W(\mathtt v)=1$, $W|_D=T$, and $D$ has a positive
$T$-grading.  The Rees total space
$\mathfrak X=\Spec\Rees_W S$ is normal and Gorenstein.  Its
parameter $\lambda$ has action weight $-1$, and the action is trivial
on $S$.  The central ring $G$ has canonical filtration shift $n$.
We also use the family $\Spec\bar S\to\A^1_t$ defined by $t=\bar v$,
whose parameter has $T$-degree zero.  Away from $t=0$ it has the graded description
\begin{equation}
 \label{boot:polynomial-family}
 \bar S[\bar v^{-1}]
 \simeq \CC[\bar v,\bar v^{-1}][u_1,\ldots,u_N],
 \qquad T(\bar v)=0,\quad T(u_i)=1.
\end{equation}
The equivariance in Equation~\eqref{boot:polynomial-family} will be
essential below.

\subsection{A slice through a point with nontrivial stabilizer}

By Corollary~\ref{filt:stabilizer}, there is a nonvertex closed point
$y_D\in\Spec D$ whose $T$-grading stabilizer is $\mu_h$, with
$h\geq2$.  Set $y=(y_D,0)\in\Spec D[\mathtt v]$, considered as
a point of the central fiber of $\mathfrak X$.  Its stabilizer for
the Rees action is the same $\mu_h$.

\begin{lemma}[The uniformizing slice and its tests]
\label{boot:slice-tests}
There is a normal Gorenstein algebraic germ $y\in Y$ of dimension
$n$, preserved by $\mu_h$, with a nonzero function
$\lambda_Y\in\fm_{Y,y}$ of character $-1$, such that $Y$ is klt
away from $\lambda_Y=0$.  Every quasi-monomial valuation $u$ centered
at $y$ gives a quasi-monomial valuation $w$ centered at $o\in C$.
Writing $\ell=u(\lambda_Y)>0$, one has
\begin{equation}
 \label{boot:basic-test}
 A_C(w)\leq N\ell+A_Y(u),
 \qquad
 \vol(w)\leq \ell^{-n}.
\end{equation}
The comparison is valid before klt at $y$ has been established.
\end{lemma}

\begin{proof}
Choose a positive-$W$-degree homogeneous function $f$ on
$\mathfrak X$ with $f(y)\neq0$.  Such a function exists because
$y_D$ is not the vertex of the positively graded ring $D$; lift a
nonvanishing homogeneous function from the central fiber.  Write
$d=W(f)>0$ and $c=f(y)$.  On the open set
$U=\{f\neq0\}\subset\mathfrak X$, take the level slice
$Y_0=\{f=c\}$.  The action map
\[
 \theta:Y_0\times\Gm\longrightarrow U,
 \qquad (q,a)\longmapsto a\cdot q,
\]
is finite étale: it is the cover obtained by adjoining a root
$a^d=f/c$ of an invertible function.  We take the germ of $Y_0$
at $y$ and denote it by $Y$.  Normality and the Gorenstein property
descend from the étale product description, and
$\dim Y=\dim\mathfrak X-1=n$.  The subgroup $\mu_h$ fixes $y$
and preserves the level slice, since $h$ divides $d$.

The restriction $\lambda_Y$ of $\lambda$ has character $-1$.
It is a nonzero divisor on $Y$: its pullback to
$Y\times\Gm$ differs from the pullback of $\lambda$ only by
the invertible group character $a^{-1}$.  It vanishes at $y$.
Since $\mathfrak X_{\lambda\neq0}\simeq C\times\Gm$ is klt,
the étale product description also shows that $Y$ is klt away
from $\lambda_Y=0$.

We record the function-field relation to $C$.  In
$K(C)(\lambda)$ the action fixes $K(C)$ and gives $\lambda$
weight $-1$, so
\[
 f=\lambda^{-d}s\quad\text{for some }s\in S,
 \qquad
 \lambda_Y^d=s/c\quad\text{on }Y.
\]
After inverting $\lambda_Y$, each component of the slice is finite
over the open set $\{s\neq0\}\subset C$.  It dominates this open
set: the equation is monic in $\lambda_Y$, and the resulting finite
algebra is flat over $S[s^{-1}]$.  Moreover $\lambda_Y$ does not
vanish identically on any component through $y$.  Thus $K(Y)$ is a
finite extension of $K(C)$, and $S$ injects into $K(Y)$.

Extend $u$ to $K(Y)(a)$ by the Gauss rule with $a$ of value zero
and residue transcendental over the residue field of $u$.
This is the extension trivial on the group factor.  Denote by $V$
its restriction through $\theta$ to $K(\mathfrak X)$.  Its center
is the generic point of the orbit through $y$, and
\begin{equation}
 \label{boot:etale-discrepancy}
 A_{\mathfrak X}(V)=A_Y(u),
 \qquad V(\lambda)=\ell.
\end{equation}
Indeed, on the product a log smooth model defining $u$ may simply
be multiplied by $\Gm$; the group parameter contributes no
discrepancy.  Pulling back log smooth models by the finite étale
chart preserves the canonical divisor and discrepancy in the
restricted-value scaling.  These assertions do not require $u$
to be $\mu_h$-invariant.

Let $w=V|_{K(C)}$.  Since elements of $S$ are invariant under the
Rees action, their pullbacks to $Y\times\Gm$ are independent of
$a$, and
\begin{equation}
 \label{boot:restriction-values}
 w(s')=u(s'|_Y)\qquad(s'\in S).
\end{equation}
Equivalently, $w$ is the restriction of $u$ along the finite
extension $K(Y)/K(C)$.  The Abhyankar equality is preserved by
finite extension and restriction, so $w$ is quasi-monomial
\cite[Proposition~3.7]{JM12}.

Write $F^pS=\{W\geq p\}$ for the second filtration.  If
$s'\in F^pS$, the Rees function $\lambda^{-p}s'$ is regular on
$\mathfrak X$, hence its restriction is regular at $y$.  Therefore
\begin{equation}
 \label{boot:filtration-domination}
 w(s')\geq p\ell.
\end{equation}
Since $F^1S$ is the vertex maximal ideal, this proves that $w$ is
centered exactly at $o$.  Applying
Lemma~\ref{filt:discrepancy} with canonical shift $Q=n$, and then
Equation~\eqref{boot:etale-discrepancy}, gives the discrepancy
inequality in Equation~\eqref{boot:basic-test}.

Finally, Equation~\eqref{boot:algebra-data} implies
\begin{equation}
 \label{boot:hilbert-cutoff}
 \dim_\CC S/F^pS=\frac{p^n}{n!}+O(p^{n-1}).
\end{equation}
For $p=\lceil k/\ell\rceil$, Equation~\eqref{boot:filtration-domination}
gives $F^pS\subseteq\fa_k(w)$.  The valuation ideal has support
only at $o$, so its colength is the same before and after
localizing $S$ at the vertex.  Taking leading coefficients in
Equation~\eqref{boot:hilbert-cutoff} proves
$\vol(w)\leq\ell^{-n}$.
\end{proof}

\begin{lemma}[The slice is klt]
\label{boot:slice-klt}
The germ $y\in Y$ is klt.
\end{lemma}

\begin{proof}
If it were not klt, there would be a prime divisor $E$ on a log
resolution of $Y$ with $A_Y(E)\leq0$ and center containing $y$.
By Lemma~\ref{boot:slice-tests}, its center is contained in
$\lambda_Y=0$, so $a=\ord_E(\lambda_Y)>0$.  Choose a log
resolution also resolving this function and the ideal of $y$,
and choose a closed point $q$ of $E$ over $y$.  In local regular
parameters at $q$, give a parameter defining $E$ cost $1$ and the
other parameters positive costs tending to zero.  The resulting
quasi-monomial valuations $u_\varepsilon$, centered exactly at
$y$, satisfy
\[
 A_Y(u_\varepsilon)=A_Y(E)+O(\varepsilon),
 \qquad
 \ell_\varepsilon=u_\varepsilon(\lambda_Y)
                  =a+O(\varepsilon).
\]
Their restrictions $w_\varepsilon$ are centered at $o$.
Since $C$ is klt, $A_C(w_\varepsilon)>0$; hence
Equation~\eqref{boot:basic-test} either already gives a
contradiction through a nonpositive upper bound for this
discrepancy, or gives
\[
 2N^n\leq\nvol(o,C)
 \leq A_C(w_\varepsilon)^n\vol(w_\varepsilon)
 \leq
 \left(N+
 \frac{A_Y(u_\varepsilon)}{\ell_\varepsilon}\right)^n.
\]
The quantity inside the last power is positive whenever this
chain is applicable, and its limit superior is at most $N$.
Letting $\varepsilon$ tend to zero contradicts $2N^n>N^n$.
\end{proof}

\begin{lemma}[Singularity of the slice]
\label{boot:slice-singular}
The germ $y\in Y$ is singular.
\end{lemma}

\begin{proof}
Suppose that $Y$ were smooth at $y$.  The étale product chart
then implies that $\mathfrak X$ is smooth at the corresponding
point.  In its regular local ring, $\lambda$ is an eigenfunction
of character $-1$ for $\mu_h$.  Lift the product parameter
$\mathtt v\in G=D[\mathtt v]$ equivariantly to a function
$\widetilde v$ of character $1$.  Its differential is nonzero:
the class of $\mathtt v$ is nonzero in the cotangent space of
$\Spec G$ at $(y_D,0)$, which is a quotient of the cotangent
space of $\mathfrak X$.

Consequently $H=\{\widetilde v=0\}\subset\mathfrak X$ is
smooth at $y$, and the germ of $\Spec D$ at $y_D$ is the
hypersurface cut out in $H$ by $\lambda|_H$.  This is a
nonzero equation: $\lambda$ and $\widetilde v$ form a regular
sequence because $\mathtt v$ is regular in $G$.  If
$\lambda|_H$ has a nonzero linear part, $D$ is smooth at
$y_D$.  Otherwise it is a minimal hypersurface equation whose
character is $-1\pmod h$, and this character is nontrivial.

We now use the other deformation $\bar S$ over
$t=\bar v$.  Its special fiber is $D$, and $t$ is invariant
under the $T$-action and therefore under $\mu_h$.  Complete its
local ring at the point $y_D$ of the special fiber.  The
completion is flat over $\CC[[t]]$.  Choose a minimal system
of eigen-coordinates of the completed local ring of $D$ and
lift them equivariantly.  Complete Nakayama gives a surjection
\[
 P=\CC[[t,x_1,\ldots,x_e]]\longrightarrow
 \widehat{\cO}_{\Spec\bar S,y_D}
\]
with kernel $I$, where $t$ is fixed and every $x_i$ is an
eigen-coordinate.  Flatness identifies $I/tI$ with the kernel
of the special-fiber presentation.

If $D$ is smooth at $y_D$, this kernel is zero and Nakayama
gives $I=0$.  In the other case it is generated by the minimal
hypersurface relation of character $-1$.  Lift that generator
to $I$ and average it to its character to obtain a relation
$\widetilde F$ of character $-1$.  Nakayama gives
$I=(\widetilde F)$.  Because the base parameter is invariant
and $-1$ is a nontrivial character, one has
\[
 \widetilde F(t,0,\ldots,0)=0.
\]
In both cases, substituting $x_1=\cdots=x_e=0$ defines a
$\mu_h$-fixed formal section over $\CC[[t]]$ through $y_D$.

Since $y_D$ is not the vertex, there is a positive-$T$-degree
homogeneous element of $D$ nonzero at $y_D$.  Lift it
homogeneously to $\bar S$.  Its value on the formal section
is a unit, so remains nonzero after passing to $\CC((t))$.
On the other hand, Equation~\eqref{boot:polynomial-family}
identifies this generic fiber, with its $\mu_h$-action, with
standard affine $N$-space: each coordinate has character $1$.
For $h\geq2$ its fixed-point scheme is exactly the origin,
where every positive-$T$-degree function vanishes.  This is
the required contradiction.
\end{proof}

\subsection{The character contribution to colength}

By Lemma~\ref{boot:slice-klt}, a minimizing valuation $u_*$
exists at $y\in Y$ and is quasi-monomial
\cite[Theorem~1.2(1)]{XZStable}.  It is invariant under
$\mu_h$: this follows from uniqueness up to scaling and
invariance of log discrepancy under automorphisms, or directly
from \cite[Theorem~1.1 and Corollary~1.2]{XZ21}.  Write
\[
 A'=A_Y(u_*),\qquad
 \ell=u_*(\lambda_Y)>0,\qquad
 \sigma=\vol(u_*)>0,
\]
and let $w$ be the valuation of $C$ constructed in
Lemma~\ref{boot:slice-tests}.

\begin{lemma}[Equal leading character lengths]
\label{boot:character-shares}
Let $R=\cO_{Y,y}$ and let $\fa_t=\{g\in R:u_*(g)\geq t\}$.
For every character $\chi\in\ZZ/h\ZZ$,
\[
 \dim_\CC(R/\fa_t)_\chi
 =\frac{\sigma}{h\,n!}t^n+o(t^n)
 \qquad(t\longrightarrow\infty).
\]
\end{lemma}

\begin{proof}
The valuation ideals are $\mu_h$-invariant, so these character
quotients are defined.  Multiplication by $\lambda_Y^j$ gives
an injection
\[
 (R/\fa_{t-j\ell})_\chi
 \lhook\joinrel\longrightarrow
 (R/\fa_t)_{\chi-j}
 \qquad(j\geq0).
\]
Indeed, the value of a product is the sum of the values, so
the kernel is exactly the indicated valuation cutoff.  For
any two characters one may choose $0\leq j\leq h-1$ connecting
them.  With $H_0=(h-1)\ell$, this gives, for each fixed
$\chi$,
\[
 \length(R/\fa_{t-H_0})
 \leq h\dim_\CC(R/\fa_t)_\chi
 \leq \length(R/\fa_{t+H_0}).
\]
The total length has asymptotic
$\sigma t^n/n!+o(t^n)$, proving the assertion.  The argument
uses the primitive character of the nonzero function
$\lambda_Y$, and requires no assumption on the tangent
representation.
\end{proof}

\begin{proposition}[The refined slice estimate]
\label{boot:refined-estimate}
For the minimizing slice valuation and its restriction,
\begin{equation}
 \label{boot:refined-volume}
 \vol(w)\leq
 \min_{0\leq x\leq1}
 \left\{\frac{x^n}{\ell^n}
                 +(1-x)^n\frac{\sigma}{h}\right\}.
\end{equation}
\end{proposition}

\begin{proof}
Set $\mathfrak a_k=\{s\in S:w(s)\geq k\}$.  For a positive
integer $p$, the colength splits exactly as
\begin{equation}
 \label{boot:exact-split}
 \dim_\CC S/\mathfrak a_k
 =\dim_\CC S/(F^pS+\mathfrak a_k)
  +\dim_\CC F^pS/(F^pS\cap\mathfrak a_k).
\end{equation}
The first term is bounded by $\dim S/F^pS$.  For the second,
restriction and division define a linear map
\[
 F^pS\longrightarrow R,
 \qquad s\longmapsto
            \lambda_Y^{-p}(s|_Y).
\]
It is regular because $\lambda^{-p}s$ belongs to the Rees
algebra.  Its image has character $p\pmod h$, since $s$ is
invariant and $\lambda_Y$ has character $-1$.  By
Equation~\eqref{boot:restriction-values}, it induces the
injection
\begin{equation}
 \label{boot:quotient-injection}
 \frac{F^pS}{F^pS\cap\mathfrak a_k}
 \lhook\joinrel\longrightarrow
 \left(R/\fa_{k-p\ell}\right)_{p\bmod h}.
\end{equation}
In particular, this step is a comparison of one character
quotient, not a comparison with the full length upstairs.

Fix $0<x<1$ and take $p=\lfloor kx/\ell\rfloor$.  Apply
Equation~\eqref{boot:hilbert-cutoff} to the first term in
Equation~\eqref{boot:exact-split}, and
Lemma~\ref{boot:character-shares} to
Equation~\eqref{boot:quotient-injection}.  There are only
finitely many characters for the fixed slice, so the error
is uniform as $p\bmod h$ varies.  Multiplication by $n!/k^n$
and passage to the limit give
\[
 \vol(w)\leq \frac{x^n}{\ell^n}
                  +(1-x)^n\frac{\sigma}{h}.
\]
Continuity extends this bound to the endpoints.  Taking its
minimum proves Equation~\eqref{boot:refined-volume}.
\end{proof}

\begin{lemma}[Optimization]
\label{boot:holder}
For the fixed cone and large-stabilizer construction,
\begin{align}
 d(o,C)
 &\leq
 \frac{\bigl((N/n)\ell+A'/n\bigr)^n}
 {\bigl(\ell^{n/N}+(h/\sigma)^{1/N}\bigr)^N}
 \label{boot:optimized}\\
 &\leq (N/n)^n+(A'/n)^n\sigma/h
 \leq \frac{p_n}{2}+\frac{M_n}{2}.
 \label{boot:bootstrap-inequality}
\end{align}
Equality in the middle, Hölder, inequality holds precisely when
\begin{equation}
 \label{boot:holder-equality}
 (A')^N\ell\,\frac{\sigma}{h}=N^N.
\end{equation}
\end{lemma}

\begin{proof}
For positive $\ell$ and $\sigma$, direct differentiation gives
\[
 \min_{0\leq x\leq1}
 \left\{\ell^{-n}x^n+\frac{\sigma}{h}(1-x)^n\right\}
 =\bigl(\ell^{n/N}+(h/\sigma)^{1/N}\bigr)^{-N}.
\]
The minimum is attained at the unique interior point
\begin{equation}
 \label{boot:optimal-split}
 x=\frac{\ell^{n/N}}
         {\ell^{n/N}+(h/\sigma)^{1/N}}.
\end{equation}
Together with the discrepancy bound in
Equation~\eqref{boot:basic-test}, this proves
Equation~\eqref{boot:optimized}.

Apply Hölder with conjugate exponents $n/N$ and $n$ to the
vectors
\[
 \bigl(\ell,(h/\sigma)^{1/n}\bigr),
 \qquad
 \bigl(N,A'(\sigma/h)^{1/n}\bigr).
\]
Their scalar product is $N\ell+A'$.  This gives
\[
 (N\ell+A')^n
 \leq
 \bigl(\ell^{n/N}+(h/\sigma)^{1/N}\bigr)^N
 \bigl(N^n+(A')^n\sigma/h\bigr),
\]
with equality exactly as in
Equation~\eqref{boot:holder-equality}.
Finally, Lemmas~\ref{boot:slice-klt} and
\ref{boot:slice-singular} show that $y\in Y$ is a singular
boundary-zero klt germ of dimension $n$.  Since $u_*$ computes
its minimum,
\[
 (A'/n)^n\sigma=d(y,Y)\leq M_n.
\]
Using $h\geq2$ proves Equation~\eqref{boot:bootstrap-inequality}.
\end{proof}

\subsection{The upper bound without a same-dimensional assumption}

\begin{theorem}[The upper bound]
\label{boot:upper}
Assume the local ODP volume theorem in all dimensions smaller
than $n$.  Then $M_n\leq p_n$.  Equivalently, every singular
boundary-zero klt $n$-dimensional germ satisfies
\[
 \nvol(x,X)\leq 2(n-1)^n.
\]
\end{theorem}

\begin{proof}
Suppose $M_n>p_n$.  Choose singular klt germs whose densities
$d_i>p_n$ tend to $M_n$.  Proposition~\ref{cone:reduction}
gives for each of them a singular Gorenstein K-polystable cone
$o_i\in C_i$, smooth off the vertex, with the same density
$d_i$.

Fix $i$ and choose the fast gradings of
Lemma~\ref{cone:approximation} for this cone.  If
$m_{\max}\leq r/N$ holds along an infinite subsequence,
Theorem~\ref{small:classification} implies that $C_i$ is
smooth or the standard ODP cone.  The first possibility
contradicts its singularity.  The second contradicts
$d_i>p_n$, since the normalized volume at an ODP is
$2(n-1)^n$ by Lemma~\ref{end:odp-value}.
Thus the large-stabilizer condition holds for all sufficiently
far terms.  Applying Lemma~\ref{boot:holder} to one such
term gives
\[
 d_i\leq p_n/2+M_n/2.
\]
Letting $i$ tend to infinity yields
$M_n\leq p_n/2+M_n/2$, contrary to $M_n>p_n$.

All choices of sufficiently fast gradings in this argument
are made separately for each fixed cone $C_i$.  Neither an
attained supremum nor approximation estimates uniform over
different cones are required.  The slice bound used only
the definition of $M_n$, so the argument assumes no
same-dimensional ODP estimate.
\end{proof}

\subsection{Equality, a strict count, and the maximal stabilizer}

Suppose now that the fixed cone satisfies $d(o,C)=p_n$.
Theorem~\ref{boot:upper} is available.  For every
sufficiently fast grading in the large-stabilizer range,
take $m=m_{\max}$ and make the construction above.

\begin{lemma}[Consequences of equality]
\label{boot:equality-data}
One has $h=2$, $d(y,Y)=p_n$, and $A'=N\ell$.
After rescaling $u_*$ so that $\ell=1$, the resulting
restriction $w$ satisfies
\begin{equation}
 \label{boot:equality-identities}
 A'=N,
 \qquad \sigma=2,
 \qquad A_C(w)=2N,
 \qquad \vol(w)=2^{-N},
 \qquad w=\frac{2N}{n}\wt_\xi.
\end{equation}
The optimal splitting fraction in
Equation~\eqref{boot:refined-volume} is $x=1/2$, and all
inequalities producing that estimate are equalities at the
level of leading coefficients.
\end{lemma}

\begin{proof}
The left endpoint of
Equations~\eqref{boot:optimized}--\eqref{boot:bootstrap-inequality}
is now $p_n$, and the right endpoint is at most $p_n$.
Every inequality is therefore an equality.  In particular,
\[
 h=2,
 \qquad (A')^n\sigma=2N^n.
\]
Combining this with Equation~\eqref{boot:holder-equality}
gives
\[
 N^N=(A')^N\ell\,\sigma/2
     =N^n\ell/A',
\]
so $A'=N\ell$.  Rescaling $u_*$ by $\ell^{-1}$ makes
$\ell=1$, hence $A'=N$ and $\sigma=2$.
Equation~\eqref{boot:optimal-split} then gives $x=1/2$ and
the refined volume bound becomes $\vol(w)\leq2^{-N}$.
The discrepancy bound becomes $A_C(w)\leq2N$.

Since $w$ is an admissible centered valuation,
\[
 2N^n=\nvol(o,C)
 \leq A_C(w)^n\vol(w)
 \leq (2N)^n2^{-N}=2N^n.
\]
Both factors are positive.  Thus both displayed upper
bounds are attained, and $w$ computes the normalized
volume minimum.  Uniqueness up to scaling
\cite[Theorem~1.1]{XZ21}, together with
$A_C(\wt_\xi)=n$, gives the last identity in
Equation~\eqref{boot:equality-identities}.
\end{proof}

\begin{lemma}[The distinguished section cannot have value greater than two]
\label{boot:strict-count}
In the normalization of Lemma~\ref{boot:equality-data},
the section $v\in S_m$ used to define the two filtrations
satisfies $w(v)\leq2$.
\end{lemma}

\begin{proof}
Suppose $w(v)>2$.  Choose a fixed real number $t$ with
\[
 \frac1{w(v)}<t<\frac12,
\]
and set $a_k=\lceil tk\rceil$ and $p_k=\lfloor k/2\rfloor$.
For large $k$ one has $p_k-a_k>0$.  Use the vector-space
basis $v^b e_i$ supplied by
Lemma~\ref{filt:first-ring}; it is adapted to $W$, with
$W(v^b e_i)=b+T_i$.  Let $E_k$ be the span of its basis
vectors with level less than $p_k-a_k$.

Multiplication by $v^{a_k}$ maps $E_k$ injectively into
$S/F^{p_k}S$.  To see injectivity modulo this filtration,
the initial form of $v$ is the regular polynomial variable
$\mathtt v$ in $G=D[\mathtt v]$.  Hence multiplication by
$v^{a_k}$ raises every nonzero initial level by exactly
$a_k$ and kills no initial form.  Moreover, every element
of $v^{a_k}E_k$ has $w$-value greater than $k$: $w$ is
nonnegative on $S$ and $a_kw(v)>k$.

Thus the kernel of
\[
 S/F^{p_k}S\longrightarrow S/(F^{p_k}S+\mathfrak a_k)
\]
has dimension at least
\[
 \dim E_k
 =\dim S/F^{p_k-a_k}S
 =\frac{(1/2-t)^n}{n!}k^n+o(k^n).
\]
This is a strict positive leading-order loss in the first
term of the exact splitting
Equation~\eqref{boot:exact-split}.  The second term still
satisfies the character estimate in
Equation~\eqref{boot:quotient-injection}.  At $x=1/2$,
$\ell=1$, and $\sigma/h=1$, these estimates give
\[
 \vol(w)
 \leq 2^{-n}+2^{-n}-(1/2-t)^n
 <2^{-N},
\]
contradicting Lemma~\ref{boot:equality-data}.
Only the single exponent $a_k$ is used at each cutoff, so
there is no overlap between collections of powers to count.
\end{proof}

\begin{lemma}[Rounding the maximal stabilizer]
\label{boot:rounding}
Along the fast gradings of the fixed equality cone, the
large-stabilizer construction implies
$Nm_{\max}\leq r$ for every sufficiently far term.
\end{lemma}

\begin{proof}
Choose fixed torus eigen-generators $g_1,\ldots,g_q$ of
$S$, with characters $\chi_1,\ldots,\chi_q$.  Write
\[
 \zeta=\frac rn\xi+\varepsilon,
 \qquad \|\varepsilon\|\longrightarrow0.
\]
The $\xi$-weights of the generators are positive.
Consequently there are fixed constants $c_1,c_2>0$ such
that their $\zeta$-degrees lie between $c_1r$ and $c_2r$
for all sufficiently far gradings.  Also
$m=m_{\max}\leq\max_i\deg_\zeta g_i=O(r)$, by
Lemma~\ref{cone:approximation}.

The section $v\in S_m$ has a polynomial lift in these
generators homogeneous for $\zeta$.  Every monomial in
such a lift uses a bounded number of generators, because
its degree is $m=O(r)$ and each generator degree is at
least $c_1r$.  Its possible torus characters therefore
belong to a fixed finite set, independent of the grading
and of the chosen section $v$.  If $\chi$ is any character
occurring nontrivially in $v$, then
\[
 m=\langle\zeta,\chi\rangle
   =\frac rn\langle\xi,\chi\rangle
       +\langle\varepsilon,\chi\rangle,
\]
and hence, uniformly for these characters,
\[
 \langle\xi,\chi\rangle
 =\frac{nm}{r}+O(\|\varepsilon\|/r).
\]
Taking the minimum over the nonzero character components
of $v$ is precisely $\wt_\xi(v)$; the same estimate
therefore holds even when $v$ is not itself a torus
eigenfunction.  Lemma~\ref{boot:equality-data} gives
\begin{equation}
 \label{boot:changing-section-error}
 w(v)=\frac{2Nm}{r}
          +O(\|\varepsilon\|/r)
       =\frac{2Nm}{r}+o(1/r).
\end{equation}
By Lemma~\ref{boot:strict-count}, $w(v)\leq2$.
Thus $Nm-r\leq O(\|\varepsilon\|)=o(1)$.
The left side is an integer, so it is nonpositive for
all sufficiently far terms.  On an exact rational ray
the error is zero, and the same conclusion is immediate.
All constants here concern the fixed cone $C$ only.
\end{proof}

\begin{theorem}[The equality cone]
\label{boot:equality-cone}
If a singular cone obtained in
Proposition~\ref{cone:reduction} satisfies $d(o,C)=p_n$,
then it is the standard ordinary double point cone.
\end{theorem}

\begin{proof}
Take the fast primitive gradings of
Lemma~\ref{cone:approximation}.  A term either
already has $Nm_{\max}\leq r$, or lies in the
large-stabilizer range.  In the latter range, all
constructions above are available for sufficiently far
terms, and Lemma~\ref{boot:rounding} again gives
$Nm_{\max}\leq r$.  The entire sequence therefore
eventually satisfies the small-stabilizer hypothesis.
Theorem~\ref{small:classification} identifies $C$ with
the smooth affine cone or the standard quadratic cone.
Since $o\in C$ is singular, only the ordinary double
point remains.
\end{proof}
